\section{系数迹比较}\label{sec:c-comparison}

本节记录 Johnson 替换领圈的 C 论证。

\subsection{乘积 Hattori 等价}

紧凑词给出字母计数
\[
p_y=42+2=44,\qquad p_z=269+2=271.
\]
在循环 $m_2$ 词中，42 个 $y/Y$ 事件每一个都紧接其后跟一个不同的 $z/Z$ 事件；$r_{xy}$ 的两个 $y/Y$ 事件亦然。见证列出全部 44 对。
这些对在 P0 控制股沿其记录法向的 44 条乘积带中实现，P0 立方体外另置 $227$ 个模型 $z$-圆。范围是 Johnson 替换领圈，而非 $\partial W_2$ 中的嵌入切割。

令 $\mathcal C$ 为对象 $T:P_{86}\to P_{88}$ 的预分次有理 tangle 范畴，$F$ 为这些矩形构成的带框西-东 tangle。MWW 一柄定理将切割系数识别为
\[
M_R(T,T')=\KhR_2(R\cup T'\cup\overline T)
\]
及其所示移位与粘合作用
\cite[Theorem~4.7]{ManolescuWalkerWedrich2023}。枢轴 tangle 对偶与 $\Q$ 上 K\"unneth 给出此替换领圈上的分次同构
\begin{equation}\label{eq:HattoriActual}
H_{T,T'}:M_R(T,T')\xrightarrow{\cong}
\operatorname{Hom}_{\mathcal C}(FT,FT')\{-44\}
\otimes A^{\otimes227},
\qquad A=\Q[X]/(X^2).
\end{equation}
这些映射构成替换领圈上的系数双模同构。先在源 link 上粘合 $f:S\to T$ 再施加乘积同痕，与 rel 边界同痕于先施加同痕再以 $Ff$ 预复合。右作用是对后复合的同样论证。故两个作用方块都是同一同时曲面的粘合结合性。在链层面两个方块为
\[
\begin{array}{ccc}
C_R(T,T')&\xrightarrow{\ H_{T,T'}\ }&
C(FT,FT')\otimes C(S^1)^{\otimes227}\\
\big\downarrow {L_f}&&\big\downarrow {L_{Ff}\otimes1}\\
C_R(S,T')&\xrightarrow{\ H_{S,T'}\ }&
C(FS,FT')\otimes C(S^1)^{\otimes227},
\end{array}
\qquad
\begin{array}{ccc}
C_R(T,T')&\xrightarrow{\ H_{T,T'}\ }&
C(FT,FT')\otimes C(S^1)^{\otimes227}\\
\big\downarrow {R_g}&&\big\downarrow {R_{Fg}\otimes1}\\
C_R(T,U)&\xrightarrow{\ H_{T,U}\ }&
C(FT,FU)\otimes C(S^1)^{\otimes227}.
\end{array}
\tag{17}
\]
此处 $C_R$ 与 $C$ 为 strict Blanchet--Khovanov foam 复形。每个方块中两个复合都是在交换两个不相交领圈后得到同一粘合 foam。BHPW strict 函子性消去 projective 符号，取同调得 \eqref{eq:HattoriActual} 及 Johnson 替换领圈的两个 MWW 作用方块。方程~\eqref{eq:HattoriActual} 不是 $\partial W_2$ 中嵌入切割的链层面复形。

普通商陈述（含两个循环关系 span）在
形式化发展中核检查。\footnote{见
\texttt{Smooth4PC/RepresentableCoefficient.lean}。}

\begin{lemma}[实际系数双模]\label{lem:C1}
对 P0 替换，\eqref{eq:HattoriActual} 是实际 MWW 系数双模（含两个粘合作用）的自然同构。此外，下述映射 \(\operatorname{Sh}_q\) 是该双模的量子 trace，而非维数匹配的替代。
\end{lemma}

\begin{remark}[引理~\ref{lem:C1} 的状态]
\Open。已确立：由 P0 控制股及其乘积法向 $z$-侧得到的 44 个模型矩形，给出 Johnson $y$-后-$z$ 配对与 $227$ 个模型 $z$-圆，以及第~\ref{sec:exact-algebra} 节的可表 $HH_0$ 约化。未确立：对实际切割 tangle $R$ 的同痕
\[
R\cup\overline T\cup T'\;\cong\;\overline{FT}\circ FT'\sqcup 227\text{ 个圆}
\]
（固定边界、detector 球与粘合支撑），从而链层面作用方块 (17) 亦未确立。没有该同痕，\eqref{eq:HattoriActual} 只是模型同构，下述映射 $\operatorname{Sh}_q$ 尚未被识别为实际系数双模的量子 trace。
\end{remark}

\subsection{量子 trace 与端点映射}

令 $R_q=\Q[q,q^{-1}]$。预分次循环关系为
\[
L_f(m)=q^{|f|}R_f(m).
\]
由引理~\ref{lem:C1}，\eqref{eq:HattoriActual} 与 counit 是 Johnson 替换领圈上的齐次系数态射。BPW 典范竖直-水平函子与 BHPW strict tangle/foam 函子给出
\begin{equation}\label{eq:actualShadow}
\operatorname{Sh}_q:q\operatorname{Tr}(\mathcal C;M_R)
\longrightarrow\operatorname{Hom}_{R_q}(E_{86},E_{88}),
\end{equation}
在 flat-base qHH$_0$/Chern 识别之后
\cite{BeliakovaPutyraWehrli2019,BeliakovaHogancampPutyraWehrli2023}。此处
\[
E_{86}=(V^{\otimes86})_{\lambda=86},\qquad
E_{88}=(V^{\otimes88})_{\lambda=86};
\]
前者秩为一，后者秩为 88。tangle 映射在限制到这些权空间之前是 $U_q(\mathfrak{sl}_2)$ 交织子。

基变换 $q\mapsto1+h$ 经有理局部化与 Noether 局部环完备分解，故为 flat。trace 余核对 $h$ 约化仅将 $q^{|f|}$ 换为 1，给出普通 MWW 系数 $HH_0$；张量积右正合性对此最后断言已足够。

\subsection{所选类与除法检测器}

令 $U:P_{86}\to P_{88}$ 为 P0 端点序中的所选单杯 tangle，并令 $T_1=F^{-1}U$。定义
\[
v_T=H_{T_1,T_1}^{-1}
  (\operatorname{Id}_U\otimes X^{\otimes227}).
\]
圆 counit 给出
$\operatorname{Sh}_h([v_T])=u_h$。其绝对次数为
\[
-44+227+315-4=494.
\]

相邻单缺陷基向量上归一化检查的 R-矩阵为
\[
\begin{pmatrix}1-q^{-2}&q^{-1}\\q^{-1}&0\end{pmatrix}.
\]
经位置基变换及 $t=q^{-2}$ 后，它是公开正 Burau 块；负块为其逆。物理缆具有混合定向。BPW 定向缆化公式将其在 writhe 依赖的分次移位与枢轴基映射下等同于全正模型。完整公开词 writhe 为零，故全局移位消失。任意次数零置换或枢轴坐标图 $P$ 同时传输算子、杯与盖：
\[
W\mapsto PWP^{-1},\qquad u\mapsto Pu,\qquad
\ell\mapsto\ell P^{-1}.
\]
矩阵系数不变；我们在该共同传输后选取公开端点标签。常数项由第~\ref{sec:exact-algebra} 节的单一领圈端点约定固定：所选 cup 连接绑定到领圈位置 $2$ 与 $87$ 的物理端点，BPW 杯/盖公式给出
\[
u_0=e_2-e_{87},\qquad
\ell_0=e_{87}^*-e_2^*,
\]
即 \eqref{eq:u} 与 \eqref{eq:ell}，不留任何未定符号。所有省略因子均为可逆 $q$-单项式或正阶 $h$ 修正，由下文引理~\ref{lem:filtered-cubic}，它们不影响 divided cubic。

在 $\Q[[h]]$ 上定义
\[
D_h=\ell_h(\rho_h(W)-I)\operatorname{Sh}_h.
\]
这在量子系数 trace 上良定义，因为 \eqref{eq:actualShadow} 已施加量子循环性。精确 Magnus 计算与 Darn\'e 的 Andreadakis 定理给出 $W\in\Gamma_3(P_{44})$ \cite[Theorem~6.2]{Darne2021}。因每个纯辫生成元作用为 $I+O(h)$，引理~\ref{lem:commutator-order} 给出单向蕴含
\[
W\in\Gamma_3(P_{44})\ \Longrightarrow\
\rho_h(W)-I\in h^3\operatorname{End}(E_{88});
\]
不断言其逆，亦不需要。独立于 Andreadakis 定理，对 $\rho_h(W)-I$ 全部 $7{,}744$ 个条目的直接计算验证了右侧的隶属关系。
故 $D_h$ 取值于 $h^3\Q[[h]]$。除以 $h^3$ 唯一，对 $h$ 约化给出与提升无关的普通泛函
\[
D_3:HH_0(\mathcal C;M_R)\longrightarrow\Q.
\]
杯、盖与基变换中的正阶修正不改变三次系数（引理~\ref{lem:filtered-cubic}）。

\begin{lemma}[滤过三次配对]\label{lem:filtered-cubic}
设 $A(h)=\sum_{j\ge0}A_jh^j$ 为 $E_{88}$ 的自同态级数且 $A_0=A_1=A_2=0$，$u(h)=\sum_k u_kh^k$ 与 $\ell(h)=\sum_i\ell_ih^i$ 为向量级数与行级数。则
\[
[h^3]\,\ell(h)A(h)u(h)=\ell_0A_3u_0,
\]
且以常数自同构 $P$ 将 $(A,u,\ell)$ 换为 $(PAP^{-1},Pu,\ell P^{-1})$ 时该系数不变。
\end{lemma}

\begin{proof}
$[h^3]\ell(h)A(h)u(h)=\sum_{i+j+k=3}\ell_iA_ju_k$，其中 $j\le2$ 的项均为零。同步传输下每一项不变。
\end{proof}

四个混合指标 Burau 符号变体
\[
-53824,\quad -53760,\quad -59008,\quad -59072
\]
非零，但它们不是冻结三次项：每一个都混用了两套端点坐标系。冻结公开收据与 Lean 在单一领圈约定下给出
\begin{equation}\label{eq:selectedD3}
[h^3]\,\ell_h(\rho_h(W)-I)u_h=2624\ne0.
\end{equation}
方程~\eqref{eq:selectedD3} 是关于冻结 Burau 计算的陈述；只有当引理~\ref{lem:C1} 将 $[v_T]$ 与实际系数迹类等同（该引理为 \Open）时它才成为 $D_3([v_T])=2624$。非零并不闭合 C 或 S。

\subsection{完整二柄余锥}

\begin{lemma}[全缆常数项]\label{lem:C2}
对每个有限缆级及每个保号缆 braid \(b\)，端点算子的常数项等于其带符号置换。将该置换标准化后，剩余纯 braid 算子为 \(I+O(h)\)。
\end{lemma}

\begin{proof}
BPW 定向缆化作用由 strict BHPW tangle 函子得到。在 \(q=1\) 时，基本表示及其对偶上的 braiding 为普通对称交换，含相反定向因子的枢轴识别。故 braid 的常数项恰为其带符号端点置换的映射。保号纯 braid 置换为恒等，故常数项为 \(I\)。所有矩阵元为 \(q\) 的 Laurent 多项式，在 \(q=1\) 正则；在 \(q=1+h\) 后减去常数项得 \(h\) 的倍数。论证与缆副本数无关，故在每个有限 MWW 级成立。此陈述仅关于端点表示，不断言反平行 MWW beta 作用的未知对称群分解。
引理应用于引理~\ref{lem:C1} 之后的 Johnson 替换。
\end{proof}

仅断言 $D_3$ 在 C1 之后下降。本小节其余部分是引理~\ref{lem:C1} 之后的 C2 论证。

P0c 在重构股上提供乘积法向，从而在替换领圈中给出不相交平行级。我们首先在每一 MWW 缆态上类型化该构造。在 owner 序 $(m_2,m_3,r_{xy},r_{yz},r_{zx})$ 下，令
\[
n_y=(42,189,2,2,0),\qquad n_z=(269,1271,2,2,0).
\]
对 $r\in\mathbb N^5$，沿 $y$- 与 $z$-余核切割实际缆 $K(r,r)$ 得
\[
p_y(r)=\sum_i r_i n_{y,i},\qquad
p_z(r)=\sum_i r_i n_{z,i}.
\tag{24}
\]
由引理~\ref{lem:C1}（一旦确立），P0 矩形的平行副本将给出与 (17) 同形的链同构 $H_r$，具相应端点集。不假设 $r_{zx}$ 分量分裂；其 stabilization 关系是第~\ref{sec:local-stabilization} 节的主题。因所有副本位于同一固定管状乘积坐标图中，$H_r$ 与每个粘合作用及添加平行对可交换。

令 $s_0=e_{m_2}+e_{r_{xy}}$。若 $r\ge s_0$，令 $\Omega_r$ 为 $m_2$ 与 $r_{xy}$ 各选一个负副本与一个正副本的全部有限选择集。对 $\omega\in\Omega_r$，用 MWW 核心 counit 闭合每个未选乘积因子，用所选 44 个通道做点推，并记所得行的三次系数为 $D_{r,\omega}$。定义
\[
D_r=|\Omega_r|^{-1}\sum_{\omega\in\Omega_r}D_{r,\omega}.
\tag{25}
\]
对不在 $s_0$ 之上的态，用 once-dotted MWW 映射补缺失对，并通过拉回 (25) 定义 $D_r$。不相交支撑使这些添加的顺序无关紧要。故 $D_r$ 在 MWW 定义~3.1 的每个直和项上定义，含所选态之下的态。

对 beta，将 P0 领圈运动同时施加于所有物理副本。这是 \cite[Theorem~E]{BeliakovaPutyraWehrli2019} 的定向缆化作用，由 BHPW 严格化。常数置换重标 $\Omega_r$。放置标准化后，剩余纯 braid 由引理~\ref{lem:C2} 为 \(I+O(h)\)。算子、杯与盖的同时传输，连同 $\rho(W)-I=O(h^3)$，仅将行改变到四阶及以上。因此
\[
D_r\beta_i(b)=D_r
\tag{26}
\]
对每个 $b\in B_{r_i,r_i}$。此三次阶论证不使用反平行 braid 作用分解过对称群的断言，MWW 对此仍开放
\cite[Remark~3.3]{ManolescuWalkerWedrich2023}。

对一对添加，按新对是否被选中分裂目标放置。beta 移动交换所选新对与旧对，故 (26) 将两个子集约化到同一局部核心计算。在整个源上该计算为
\[
(\epsilon\otimes\epsilon)\Delta(1)=0,\qquad
(\epsilon\otimes\epsilon)\Delta(X)=1.
\]
对 $r_{zx}$，同样的方程须在该分量的管状邻域内成立，且不作任何分裂假设；这是第~\ref{sec:local-stabilization} 节的局部陈述。遗忘 distinguished 新对在放置集之间有常数大小纤维，(25) 中的归一化消去其基数。dotted 原始次数 $+2$ 被目标缆化移位 $-2$ 抵消。因此对每个 owner 与态，
\[
D_{r+e_i}\psi_i^{[0]}=0,\qquad
D_{r+e_i}\psi_i^{[1]}=D_r.
\tag{27}
\]
在 $s_0$ 之下第二个等式为定义；Frobenius 恒等式与不相交支撑交换证明第一个等式及所有 owner 方块可交换。

一旦 (26)--(27) 在实际缆化源上确立，连同 (17) 的系数 trace 关系，它们在整个缆化直和上证明
\[
D_3(\beta_i(b)x-x)=0,\quad
D_3(\psi_i^{[0]}x)=0,\quad
D_3(\psi_i^{[1]}x-x)=0,
\]
商 universal 性质与 MWW 定理~3.2 随即给出 $\overline D_3:\mathcal S^2_0(W_2;\Q)\to\Q$。目前 (17) 与 (26)--(27) 仅在控制领圈上验证，故该下降属于 \Open 的假设~\ref{hyp:P1}。

\subsection{无分裂假设的局部稳定化}
\label{sec:local-stabilization}

对带框分量 $K_i$ 与缆态 $r$，令 $\chi_{i,r}$ 为 $K_i$ 的 $r_i$ 个平行副本在 $K_i$ 管状邻域内的局部闭合，在预商 foam 范畴中由无点与一点球面求值定义。每个 owner（含 $r_{zx}$）的 $\psi$ 关系所需的陈述是
\[
\chi_{i,r+e_i}\circ\psi_i^{[0]}=0,\qquad
\chi_{i,r+e_i}\circ\psi_i^{[1]}=\chi_{i,r},
\]
且不假设 $K_i$ 与 link 其余部分分裂。该局部稳定化--收缩陈述为 \Open。较早文本由运输后 $r_{zx}$ 的空自由词推断 split unknot；该推断撤回，因为自由群词不决定嵌入分量。

\subsection{C 结论与次数 $494$ 方块}

MWW 一柄与二柄映射嵌入交换图
\[
\begin{array}{ccccccc}
M_R(T_1,T_1)&\xrightarrow{H_{T_1,T_1}}&
\operatorname{End}(U)\{-44\}\otimes A^{\otimes227}
&\xrightarrow{\operatorname{id}\otimes\epsilon^{\otimes227}}&
\operatorname{End}(U)&\xrightarrow{d_3}&\Q\\
\big\downarrow{\mathrm{Glue}_{1h}}&&&&&&\big\Vert\\
\mathcal S^2_0(W_1;K(s_0,s_0))\{315\}
&\xrightarrow{\ \iota_{s_0}\{-4\}\ }&
\mathcal S^{2,0}_0(W_1;K)/(\beta,\psi)
&\xrightarrow[\cong]{\ \Phi\ }&\mathcal S^2_0(W_2)
&\xrightarrow{\overline D_3}&\Q.
\end{array}
\tag{28}
\]
左竖映射及其移位为 MWW 定理~4.7，底同构为 MWW 定理~3.2。上所选元为 $\operatorname{Id}_U\otimes X^{\otimes227}$。其逐次次数为
\[
-44+227=183,\qquad183+315=498,\qquad498-4=494.
\]
在假设~\ref{hyp:P1} 下，图 (28) 将其映到类 $x_2$，在 Johnson 替换领圈上满足 $\overline D_3(x_2)=2624\ne0$。

\begin{proposition}[C 的状态]\label{thm:Cdischarge}
假设~\ref{hyp:P1} 对 Johnson 替换为 \Open。已确立：可表 $HH_0$ 约化、端点表示上的引理~\ref{lem:C2}、引理~\ref{lem:filtered-cubic}，以及冻结 Burau 值 \eqref{eq:selectedD3}。未确立：引理~\ref{lem:C1} 的实际乘积带同痕与链层面作用方块；按放置平均定义 $D_r$ 及其在实际源每个缆态上的关系 (26)--(27)；第~\ref{sec:local-stabilization} 节的局部稳定化陈述；从而 $D_3$ 通过完备二柄余锥的下降，以及所选次数 $494$ 类在所有 MWW 二柄关系之后非零。
\end{proposition}

\begin{proof}
已确立各项是附录~\ref{sec:reproducibility-extended} 所列可重放证书与 Lean 文件。其余各项需要实际切割 tangle，它不属于当前证书。
\end{proof}

\paragraph{C/S 防火墙。}
C 中得到的 $B$-external naturality 涉及系数比较中的边界配边。它不把 $\partial W_2$ 中本质球面的局部作用等同于普通 dotted 球面求值；那是引理~\ref{lem:Sendpoint} 中的单独计算。
